\chapter{Symmetric Hashing and Key Storage}
\label{ch:5.3}
\chaptertopic{Topic 5.3}

\chapterquote{The system must not require secrecy, and it must be able to fall into the hands of the enemy without inconvenience.}{Auguste Kerckhoffs, \emph{La Cryptographie Militaire} (1883)}

\section{Motivation}

Chapters~\ref{ch:5.1} and~\ref{ch:5.2} were about who is allowed to touch data. This chapter is
about protecting the data itself, cryptographically, regardless of who eventually reaches it. Three
questions organize it. When should stored data be reversible (encryption) versus permanently
one-way (hashing)? How does symmetric encryption actually work, and how are its keys kept safe over
their whole lifetime? And how do you prove not just that data is unaltered, but that it genuinely
came from whoever claims to have sent it?

\begin{apalignment}{Unit 5: Securing Applications and Data}{Topic 5.3 -- Protecting Stored Data with Cryptography}
\begin{apobjectives}
\apobjective{5.3.A}{Mitigate Risk}{Decide whether data should be stored encrypted (recoverable) or hashed (one-way) based on whether the original value must ever be retrieved, and manage a symmetric key safely across its lifecycle -- generation, storage, and rotation -- so no single compromised secret exposes the underlying data.}
\apobjective{5.3.B}{Mitigate Risk}{Use a symmetric algorithm to encrypt and recover data, handling correctly the parts that surround the cipher itself -- where the key comes from, how the initialization vector is generated, and how the ciphertext is authenticated so that tampering is detected rather than silently decrypted.}
\end{apobjectives}
\end{apalignment}

\section{Encryption versus Hashing}

Both tools turn readable input into an unreadable-looking blob of bytes, and both are indispensable
-- but they solve opposite problems, and using the wrong one is a common, costly mistake.

\begin{definition}{Encryption}
A reversible transformation of plaintext into ciphertext using a key; given the correct key, the
ciphertext decrypts back to the exact original plaintext.
\end{definition}

\begin{definition}{Hash function}
A one-way function that maps an input of any size to a fixed-size \textbf{digest}, with no
algorithm capable of reversing the process -- given only a digest, the original input cannot be
recovered, only guessed.
\end{definition}

The question that separates them is simple: do you ever need to get the original back? A payroll
file must be readable whenever payroll actually runs, so it is stored encrypted. A login password
never needs to be recovered at all -- the site only needs to know whether a submitted password
matches what was registered -- so it is stored hashed. At login, the input is hashed and the two
digests are compared; the original password is never reconstructed anywhere in the process. If
a site stored passwords encrypted instead, a database breach combined with key compromise would
expose every user's password in plaintext at once. Hashing removes that single point of failure
entirely -- there is no key to steal that would undo it.

Hash functions rely on the \textbf{avalanche effect}: two inputs differing by a single character
produce completely unrelated digests, which is what makes a hash useful for detecting even the
smallest tampering. But a naive hash of a common password produces the same digest every time it is
hashed anywhere. That is exactly what \textbf{rainbow tables} -- enormous precomputed
password-to-hash lookup tables -- exploit to crack common passwords instantly. A \textbf{salt}, a
random value appended before hashing and stored alongside the digest, defeats this: even two users
who chose the identical password end up with different digests, because their salts differ.

\begin{lstlisting}[language=Python, caption={Salted password hashing and verification}]
import hashlib, secrets

def hash_password(password, salt=None):
    """Returns (salt_hex, digest_hex). Generates a fresh salt if none given."""
    salt = salt or secrets.token_hex(16)
    digest = hashlib.sha256((salt + password).encode()).hexdigest()
    return salt, digest

def verify_password(password, salt, expected_digest):
    _, digest = hash_password(password, salt)
    return secrets.compare_digest(digest, expected_digest)
\end{lstlisting}

The SHA-2 family (SHA-256, SHA-512) is today's standard for general-purpose hashing. MD5 and SHA-1
are deprecated for any security use because both have documented \textbf{collision} weaknesses --
two genuinely different inputs producing the identical digest.\srccite{5.3}{1} That collision
weakness undermines the entire integrity guarantee a hash is meant to provide. For passwords specifically, purpose-built slow
algorithms (bcrypt, scrypt, Argon2) are preferred over a fast general-purpose hash like SHA-256
alone, because their deliberate computational cost makes brute-force guessing expensive even on
modern hardware.

\begin{misconception}
\textbf{``Encoding (Base64) is a form of encryption.''} Encoding is a lossless, keyless format
change, trivially reversed by anyone -- it provides zero confidentiality and should never be
mistaken for either encryption or hashing.
\end{misconception}

\section{Symmetric Encryption and Key Management}

\begin{definition}{Symmetric encryption}
A cryptographic system in which a single secret key performs both encryption and decryption; the
same key must therefore be kept secret from anyone who should not read the protected data.
\end{definition}

\textbf{AES (Advanced Encryption Standard)}, established by NIST in 2001, is the modern industry
standard, supporting 128-, 192-, and 256-bit keys, all currently considered secure against
brute-force attack.\srccite{5.3}{2} It replaced DES specifically because DES's 56-bit key had become
brute-forceable with modern hardware. AES is a \textbf{block cipher}: it encrypts fixed-size 128-bit blocks.
Real data is almost always larger than one block, so a \textbf{mode of operation} defines how
blocks chain together -- and that choice matters enormously for security.

\begin{example}[Why ECB mode leaks the exact shape of your data]
\textbf{Electronic Codebook (ECB)} mode encrypts each block independently with the same key --
which means two identical plaintext blocks always produce identical ciphertext blocks. Encrypt a
bitmap image in ECB mode and the image's outline remains clearly visible in the ciphertext. Large
uniform regions of the source image become large uniform regions of ciphertext, in exactly the
same shape. ECB should never be used for multi-block data regardless of how strong the underlying
cipher is.
\end{example}

\textbf{Cipher Block Chaining (CBC)} fixes this by XORing each plaintext block with the previous
ciphertext block before encrypting it, seeded by a random, non-secret \textbf{Initialization Vector
(IV)} for the first block. The IV must be unique for every encryption -- reusing one with
the same key lets an attacker detect whether two messages start with identical data. \textbf{GCM
(Galois/Counter Mode)} goes further: it turns the block cipher into a stream cipher and is an
\textbf{Authenticated Encryption with Associated Data (AEAD)} mode. That means it provides
confidentiality \emph{and} integrity verification together through an authentication tag: if an
attacker tampers with the ciphertext at all, decryption fails outright rather than silently
producing corrupted plaintext.

A cryptographic system is only as secure as how it manages its keys across their whole lifecycle.
\textbf{Generation} must use a \textbf{CSPRNG} (Cryptographically Secure Pseudo-Random Number
Generator). An ordinary language's built-in \texttt{random} module uses a predictable algorithm
and must never generate a real key; Python's \texttt{secrets} module or an OS entropy source like
\texttt{/dev/urandom} is the correct choice. \textbf{Storage} must never mean hardcoding a key in
source code or a plaintext configuration file, since either leaks the moment version control is
compromised. Dedicated \textbf{secrets managers} or \textbf{Hardware Security Modules (HSMs)} are the
correct alternative -- both perform cryptographic operations internally, so a key never leaves the
hardware in plaintext.

\begin{lstlisting}[language=Python, caption={Envelope encryption: wrapping a data key with a key-encryption key}]
def wrap_dek(dek_bytes, kek_bytes):
    """Simplified XOR-based key wrapping for illustration.
    dek_bytes, kek_bytes: equal-length bytes objects."""
    return bytes(a ^ b for a, b in zip(dek_bytes, kek_bytes))

def unwrap_dek(wrapped_dek, kek_bytes):
    return wrap_dek(wrapped_dek, kek_bytes)
\end{lstlisting}

Large systems manage many keys through \textbf{envelope encryption}: a \textbf{Data Encryption Key
(DEK)} encrypts the actual data, and a \textbf{Key Encryption Key (KEK)} encrypts (wraps) that DEK
in turn. The wrapped DEK sits next to the encrypted data itself, while the KEK stays securely inside
an HSM or secrets manager. To decrypt, the system sends the wrapped DEK to be unwrapped by the
KEK, retrieves the plaintext DEK, and only then decrypts the actual data. Keys are \textbf{rotated}
on a schedule to limit how much data any single compromised key would expose. When a key is
finally retired, it is securely destroyed -- \textbf{cryptographic erasure} -- so that without it,
the ciphertext it protected becomes permanently unrecoverable by design.

\begin{figure}[h]
\centering
\begin{tikzpicture}[
  box/.style={draw=kickstartermaroon, thick, rounded corners, minimum width=3.1cm, minimum height=0.9cm, align=center, font=\scriptsize},
]
\node[box, fill=kickstartermaroon!25] (kek) at (0,1.6) {KEK\\stays inside HSM /\\secrets manager};
\node[box, fill=kickstartercream] (dek) at (0,0) {DEK\\(wrapped by KEK)};
\node[box] (data) at (4.4,0) {Encrypted data\\(encrypted by DEK)};
\draw[-{Latex[length=2mm]}, kickstartermaroon, thick] (kek) -- (dek) node[midway, right, font=\tiny] {wraps};
\draw[-{Latex[length=2mm]}] (dek) -- (data) node[midway, above, font=\tiny] {encrypts};
\node[font=\scriptsize\sffamily, gray, align=center] at (0,-1.1) {wrapped DEK sits next to the\\encrypted data it protects};
\end{tikzpicture}
\caption{Envelope encryption: the KEK never leaves the HSM, and only the much smaller DEK needs to
be unwrapped to decrypt the actual data -- rotating the KEK never requires re-encrypting the data
itself, only re-wrapping the DEK.}
\label{fig:5.3-envelope}
\end{figure}

\begin{example}[LastPass, 2022: When the Key Management Story Falls Apart]
Attackers breached the password manager LastPass in two stages between August and October 2022,
fully disclosed only in stages through December 2022. It became a widely cited case study in what
happens when key-management assumptions do not hold up end to end.\srccite{5.3}{3} In the first stage,
attackers compromised a software engineer's corporate laptop and used it to steal source code and
technical secrets. In the second stage, using information from the first, attackers targeted a
senior DevOps engineer, ultimately capturing credentials that gave access to cloud storage backups
containing customer vault data. Those vaults were encrypted, meant to be readable only by each
customer, using a key derived from that customer's own master password -- which LastPass itself
never stored. For most customers, that design held: without the master password, the stolen vaults
were not directly readable. But LastPass's own later disclosure revealed two compounding
weaknesses. Some vaults had been created years earlier under a much lower password-based
key-derivation iteration count than the company's current standard. Older accounts had not always
been migrated forward, leaving them meaningfully cheaper for an attacker to brute-force offline.
And within each vault, some fields -- website URLs among them -- had historically been stored
unencrypted for functional reasons, exposing which sites a customer had accounts on even without
cracking anything. The lesson generalizes past this one company. Envelope encryption and strong
algorithms only deliver the protection this chapter describes when every parameter in the chain is
actually current for every record -- not just the ones created after a policy improved. That chain
includes the key-derivation iteration count, which fields are actually encrypted, and whether older
data was ever upgraded to current standards. A cryptographic design can be correct on paper and
still fail in exactly the records an organization forgot to re-encrypt.
\end{example}

\begin{misconception}
\textbf{``AES-ECB is secure because AES is strong.''} Even with a maximally strong cipher, ECB
mode encrypts identical plaintext blocks into identical ciphertext blocks, exposing data patterns
regardless of how unbreakable the underlying AES computation itself is.
\end{misconception}

\section{Integrity and Authenticity}

A plain hash lets you detect whether data changed by accident -- a bad disk sector, a dropped byte
in transit -- but it does nothing against an intelligent attacker. Hash functions are public
and keyless: anyone, including the attacker, can compute SHA-256 over any input they choose. An
attacker who intercepts a message and its hash, edits the message, and recomputes the hash over the
tampered content produces a digest that matches perfectly. The check passes cleanly, because the
attacker computed it correctly over content they altered.

\begin{definition}{MAC (Message Authentication Code)}
A keyed hash that proves both \textbf{integrity} (the data was not altered) and
\textbf{authenticity} (the data genuinely came from whoever holds the shared secret key). A plain
hash, by contrast, can only ever demonstrate the first property -- and only over a channel an
attacker cannot also tamper with.
\end{definition}

The obvious first attempt at building one, \texttt{H(key || message)}, has a real vulnerability in
any hash built on the Merkle-Damgard construction, which includes MD5, SHA-1, and SHA-256. These
functions expose their entire internal state as the output digest. That means an attacker who
knows \texttt{H(key || message)} can compute a valid hash for \texttt{key || message || padding ||
attacker\_data} without ever learning the key -- a \textbf{length-extension attack}. The digest
already tells the attacker exactly what internal state the function was in after processing the
original input, letting them resume hashing from there.

\begin{definition}{HMAC}
The NIST FIPS 198-1 standardized keyed-hash construction,\srccite{5.3}{4}
$\mathrm{HMAC}(K, m) = H\big((K \oplus \mathrm{opad}) \,\|\, H((K \oplus \mathrm{ipad}) \,\|\, m)\big)$,
which nests the hash function twice around two independently derived keyed prefixes so that no
exposed intermediate state is ever available to extend.
\end{definition}

Because the final output is a hash \emph{of} a hash, with the key folded in on both passes, an
attacker who only ever sees the final tag has no usable internal state to extend. They cannot
produce a valid tag for any message without the key itself. The nesting is what makes HMAC provably
resistant to length extension, even while built directly on a vulnerable Merkle-Damgard hash like
SHA-256.

\begin{lstlisting}[language=Python, caption={Building and verifying an HMAC tag by hand}]
import hashlib

def hmac_sha256(key, message, block_size=64):
    key = key if len(key) <= block_size else hashlib.sha256(key).digest()
    key = key.ljust(block_size, b"\x00")
    ipad = bytes(b ^ 0x36 for b in key)
    opad = bytes(b ^ 0x5c for b in key)
    inner = hashlib.sha256(ipad + message).digest()
    return hashlib.sha256(opad + inner).hexdigest()

def verify_hmac(key, message, tag):
    import hmac as hmac_lib
    return hmac_lib.compare_digest(hmac_sha256(key, message), tag)
\end{lstlisting}

HMAC shows up constantly in practice. In API request signing, a client and server share a secret,
and every request carries a tag over its contents so the server can reject anything tampered with
or forged. In webhook verification, a SaaS platform signs a payload with HMAC-SHA256 so the
receiving server can confirm the request really came from that platform. And JWT tokens using the
\texttt{HS256} algorithm have a server-side secret sign the header and payload, so tampering with
a token's own claims is detectable. In every case, HMAC answers the same question: did this data
arrive unmodified, from someone who actually holds the shared secret? That is a question a plain
hash structurally cannot answer, once an attacker is capable of rewriting both the message and its
accompanying digest.

\begin{misconception}
\textbf{``A matching hash proves the message wasn't tampered with.''} Only if the attacker could
not also recompute the hash after tampering. A plain, keyless hash gives no such guarantee the
moment an attacker can rewrite both the message and its accompanying digest together.
\end{misconception}

\section*{Review Questions}
\begin{reviewquestions}
\item A company needs to run payroll, which requires reading employee salary figures, and also
needs to verify user login passwords. Explain which cryptographic control belongs to each task and
justify the choice using the ``do you need the original back'' test from this chapter.
\item Explain why encrypting a bitmap image in ECB mode leaves the image's outline visible in the
ciphertext, and identify the specific property of ECB mode responsible for this leak.
\item An organization hardcodes an AES key directly in its application source code, committed to a
private Git repository. Explain why this is considered unsafe key storage even though the
repository is private, and describe the correct alternative.
\item An attacker intercepts a message and its SHA-256 hash, edits the message, and recomputes a
matching hash over the tampered content. Explain why the recipient's integrity check passes despite
the tampering, and explain what HMAC adds that would have caught it.
\item Explain why \texttt{H(key || message)} is vulnerable to a length-extension attack on
SHA-256, and explain specifically how HMAC's double-nested construction closes that gap.
\item In the LastPass breach, most customer vaults remained unreadable to the attackers because
the decryption key was derived from a master password LastPass never stored. Explain why older
vaults created under a lower key-derivation iteration count were still meaningfully more exposed
than newer ones, even though both used the identical envelope-encryption design in principle.
\item Using the envelope-encryption diagram, explain why an organization can rotate its KEK on a
regular schedule without re-encrypting its entire dataset of stored files. Identify exactly
which artifact actually changes when a KEK rotation happens.
\end{reviewquestions}

\begin{answerkey}
\item Payroll requires the original salary figures to be recoverable whenever payroll actually
runs, so it is stored encrypted -- the ``do you need the original back'' answer is yes. Login
passwords never need to be recovered; the system only needs to confirm a submitted password matches
what was registered, so passwords are stored hashed. The original is never reconstructed
anywhere in the verification process.
\answerlabel{5.3.A}
\item ECB mode encrypts each fixed-size block independently under the same key, so two identical
plaintext blocks always produce identical ciphertext blocks. Large uniform regions of a bitmap
image consist of many identical or near-identical blocks, which map to identical ciphertext blocks
in exactly the same arrangement, so the image's outline survives encryption intact. The specific
property responsible is ECB's lack of chaining or randomization between blocks.
\answerlabel{5.3.B}
\item A private repository is not a secrecy boundary that survives normal operations: any developer
with repository access, any CI system that clones it, and anyone who obtains a compromised
credential or a leaked clone can read the key. The key also remains exposed in the repository's full
history even after later removal. The correct alternative is a dedicated secrets manager or
Hardware Security Module, which stores the key outside source control. In the HSM case, it also
performs cryptographic operations internally, so the key never leaves the hardware in plaintext at
all.
\answerlabel{5.3.A}
\item The check passes because SHA-256 is a public, keyless function: the attacker can compute a
correct hash over any content they choose, including their own tampered message. The recomputed
hash matches the tampered content perfectly. HMAC adds a shared secret key folded into the
computation itself, so an attacker without that key cannot produce a valid tag for their tampered
message, even though they can still compute an ordinary SHA-256 hash of it.
\answerlabel{5.3.B}
\item \texttt{H(key || message)} is vulnerable because Merkle-Damgard hash functions expose their
entire internal state as the output digest. An attacker who knows \texttt{H(key || message)}
already knows exactly what internal state the function was in after processing the original input.
They can resume hashing from there, computing a valid hash over \texttt{key || message || padding ||
attacker\_data} without ever learning the key. HMAC closes this by nesting the hash function twice
around two independently keyed prefixes, so the final tag is a hash of a hash with the key mixed in
on both passes. An attacker who sees only the final tag has no exposed intermediate state usable
as a starting point for extension.
\answerlabel{5.3.B}
\item A key-derivation function's iteration count directly controls how computationally expensive
each single guess is for an attacker performing offline brute-force against the master password.
That is the same deliberate-slowness principle bcrypt and Argon2 apply to password hashing more
generally. A vault created under a much lower iteration count requires far less computation per
guess, so the identical envelope-encryption design and identical algorithm still yield a
dramatically cheaper attack surface for older vaults. The design's security depends on a tunable
parameter that was never uniformly current across every account, not on the algorithm choice alone.
\answerlabel{5.3.A}
\item Because the KEK only wraps the (much smaller) DEK rather than touching the actual data
directly, rotating the KEK requires only unwrapping every existing DEK with the old KEK and
re-wrapping it with the new one. That is a small, fast operation compared to re-encrypting the
underlying files themselves, which can be arbitrarily large and numerous. The artifact that
actually changes during a KEK rotation is the wrapped DEK stored alongside each piece of encrypted
data. The encrypted data itself, and the DEK's own plaintext value, never need to be touched at
all.
\answerlabel{5.3.B}
\end{answerkey}

\section*{Sources}
\begin{enumerate}[label=\lbrack\arabic*\rbrack, leftmargin=2.2em]
\item \srcitem{5.3}{1}National Institute of Standards and Technology, ``NIST Transitioning Away
from SHA-1 for All Applications,'' NIST Computer Security Resource Center, December 15, 2022.
\item \srcitem{5.3}{2}National Institute of Standards and Technology, Federal Information Processing
Standards Publication 197, \emph{Advanced Encryption Standard (AES)}, November 26, 2001.
\item \srcitem{5.3}{3}LastPass, ``Security Incident Update \& Recommended Actions,'' The LastPass
Blog, March 1, 2023 (final incident report, consolidating the August--December 2022 disclosures).
\item \srcitem{5.3}{4}National Institute of Standards and Technology, Federal Information Processing
Standards Publication 198-1, \emph{The Keyed-Hash Message Authentication Code (HMAC)}, July 16,
2008.
\end{enumerate}
